\problem{Магнитные цепи и трансформатор}{10.0}

Закон Ома известен всем из школьной физики. В этой задаче, с точки зрения магнетизма, Вы ознакомитесь с понятием о \textit{магнитных цепях}, для которых применимы аналогичные электрическим цепям законы, и посредством них проанализируете некоторые применяемые в электротехнике устройства.

\pic5r{0pt}{0.2}{figures/3.1.png}

Наряду с известным Вами вектором магнитной индукции $\vec B$, полезно ввести так называемый \textit{вектор напряжённости магнитного поля} $\vec H$. Теорема Ампера для напряжённости магнитного поля утверждает, что циркуляция вектора $\vec H$ по замкнутому контуру $C$ в точности равна \textit{суммарной} силе тока проводимости, протекающей сквозь контур $C$:
$$
\oint_C\vec H\cdot d\vec l = I.
$$
В этой задаче полагайте, что все магнетики подчиняются линейному закону $\vec B = \mu\mu_0\vec H$, где $\mu$ --- магнитная проницаемость вещества, а $\mu_0=4\pi\cdot10{-7}~\unit{\H\per\m}$ --- магнитная постоянная. Во всей задаче в качестве магнетика будет выступать железо, для которого проницаемость $\mu=2000$ постоянна, а явлением ферромагнитного насыщения и гистерезиса следует пренебречь.

\part{Разминка}

\pic4r{0pt}{0.171}{figures/3.2.png}[Соленоид к пунктам \taskref{1}--\taskref{3}]

Рассмотрим цилиндрический сердечник радиуса $r=\qty{1}{\cm}$ и длины $l=\qty{50}{\cm}$, на который вдоль всей длины равномерно намотали $N=200$ витков провода. По проводу течёт постоянный ток силой $I=\qty{1}{\A}$.
\begin{task}
    \item Вычислите модуль магнитной индукции $B$ в центре соленоида, пренебрегая краевыми эффектами
\end{task}

\begin{task}
    \item Магнитный поток $\Phi$ в сердечнике можно записать в виде
    \begin{equation}\label{task:3.1}
    \Phi = \frac{NI}{R_m},
    \end{equation}
    где $R_m$ --- так называемое \textit{магнитное сопротивление}, а $\mathcal F=NI$ иногда называют \textit{магнитодвижущей силой}. Вычислите $R_m$ в рамках предыдущего пункта.
    \item Численно рассчитайте индуктивность $L$ такого соленоида.
\end{task}

\pic6r{0pt}{0.236}{figures/3.3.png}[К пункту \taskref{4}; масштабы не соблюдены]

Для сердечника более сложной и изгибистой формы, магнитный поток ``повторяет'' его форму, практически не создавая потока утечки. В этом можно убедиться в следующем пункте:
\begin{task}
    \item Рассмотрим плоскую границу воздух--железо как на Рис. \figref{2}: вектор магнитной индукции в воздухе $\vec B_0$ направлен под углом $\alpha=\ang{10}$ к нормали как показано на Рис. \figref{2}. Пусть вектор магнитного поля в железе $\vec B$ наклонён под углом $\beta$ к нормали. Чему равен $\beta$ и отношение $B/B_0$ модулей магнитных полей между двумя средами? Поверхностных токов проводимости нет.
\end{task}

\part{Сердечник с зазором}

\pic9R{0pt}{0.226}{figures/3.4.png}[К пунктам \taskref{5}--\taskref{7}; масштабы не соблюдены]

Рассмотрим сердечник в форме скругленного квадрата стороной $a=\qty{20}{\cm}$ и площадью сечения $S$. На одну сторону квадрата плотно намотаны $N$ витков провода, по которому течёт ток $I$, а с противоположной стороны разрезан зазор толщины $d$, как на Рис. \figref{3}.

\begin{task}
    \item Найдите $R_m=R_0$, всей системы, где $R_0$ определён из \eqref{task:3.1}. Покажите, что его можно записать как сумму магнитных сопротивлений сердечника и воздушного зазора.
\end{task}

\begin{task}
    \item При каком численном значении $d$ магнитные сопротивления сердечника и воздушного зазора будут равны?
\end{task}

На квадратный сердечник с теми же геометрическими размерами, как на Рис. \figref{3}, \textit{но без воздушного зазора} ($d=0$), обмотали два провода с противоположных сторон: один провод имеет $N_1$ витков и подключен к сети некоторого переменного напряжения $V(t)$; другой провод имеет $N_2$ витков и подключен к резистору сопротивлением $R$. Поток утечки отсутствует.
\begin{task}
    \item Определите коэффициент взаимной индукции $M$ двух катушек. Чему равно сопротивление цепи $R'$, видимое со стороны источника $V(t)$? \textit{Подсказка 1:} коэффициент $M$ определяется как отношение магнитного потока $\Phi_{12}$ первой катушки, сцепленной со второй, к току в первой: $M=\frac{N_2\Phi_{12}}{I_2}=\frac{N_1\Phi_{21}}{I_1}$. \textit{Подсказка 2:} для нахождения $R'$ учтите малость $R_m$ сердечника.
\end{task}

\part{Электромагнитная аналогия}

Пункты \taskref{1}--\taskref{6} наводят на мысль провести аналогию между электрическими и магнитными цепями. Впоследствии эта аналогия \textit{будет полезна для решения последующих пунктов задачи}.

\begin{task}
    \item Ниже приведена таблица соответствия электрических и магнитных величин: Вам необходимо перенести таблицу в свой рабочий лист и напротив магнитных величин записать соответствующие им электрические величины, и наоборот. \textit{Если затрудняетесь записать словесное определение, обязательно укажите стандартный символ данной величины.} Первые две строки уже заполнены за Вас; их переписывать в рабочий лист не нужно.
    \begin{center}
    \begin{tabular}{|c|c|}\hline
         \textbf{Магнитная цепь}& \textbf{Электрическая цепь} \\\hline
         Магнитодвижущая сила $NI$& Электродвижущая сила $\mathcal E$\\\hline
         Магнитное сопротивление $R_m$& Электрическое сопротивление $R$\\\hline
         Магнитный поток $\Phi$ &\\\hline
         &Напряжённость электрического поля $\vec E$ \\\hline
         Абсолютная магнитная проницаемость $\mu\mu_0$&\\\hline
    \end{tabular}
    \end{center}
    
    \item Укажите, чему равна запасённая магнитная энергия $W_m$ участка магнитной цепи сопротивлением $R_m$ при фиксированной магнитодвижущей силе $\mathcal F$.
\end{task}

\part{Дифференциальный трансформатор}

В простейшем виде дифференциальный трансформатор продемонстрирован на Рис. \figref{4}. Через сердечник в форме ``8'' размерами $2l\times l$, где $l=\qty{20}{\cm}$ и магнитным сопротивлением на единицу длины $\rho_m=\qty{4e5}{\per\H\per\m}$, обмотаны на двух противоположных концах провода с равным количеством витков $N=10$. В электротехнике, провод слева называют \textit{фазой} --- линия под напряжением --- и через него протекает переменный ток с амплитудным значением $I_1=\qty{5}{\A}$ и частотой $f=\qty{50}{Гц}$; провод справа называют \textit{нулём} --- нулевой рабочий проводник --- через который в обратном направлении течёт переменный ток с амплитудным значением $I_2=\qty{4.97}{\A}$ и частотой $f=\qty{50}{Гц}$. Считайте, что фазовый сдвиг между $I_1$ и $I_2$ отсутствует. Центр сердечника имеет вторичную обмотку с $N_0=1000$ витками, который подсоединён к последовательно соединённым резистору $R=\qty{200}{\ohm}$ и магнитным пускателем $M$. Считайте, что $M$ имеет пренебрежимо малые активное и реактивное сопротивления.
\begin{task}
    \item Определите значение амплитудной силы тока $I_M$ через магнитный пускатель $M$.
\end{task}

\cpic{0.5}{figures/3.5.png}[Дифференциальный трансформатор к пункту \taskref{10}]

\part{Магнитный пускатель}

Рассмотрим детально принцип работы магнитного пускателя, который изображён на Рис. \figref{5}. На П-образной подставке \textbf{(a)} радиуса $b=\qty{15}{\mm}$ и малой толщины закреплён железный сердечник \textbf{(c)} радиусом $a=\qty{7}{\mm}$, а через закреплённые на подставке \textbf{(a)} пружины \textbf{(e)} суммарной жёсткости $k=\qty{50}{\N\per\m}$ подсоединён, оставляя между своим внутренним торцом и верхним торцом \textbf{(c)} вертикальный воздушный зазор изначальной высоты $d=\qty{5}{\mm}$, подвижный железный сердечник внутренним радиусом чуть больше $a$ и внешним радиусом чуть меньше $b$ (т.е. попарные зазоры между подставкой \textbf{(a)} и сердечниками \textbf{(b)} и \textbf{(c)} гораздо меньше $d$). Длины сердечников гораздо больше $b$. Вокруг подставки плотно обмотаны витки \textbf{(d)} провода, который подключен к внешней цепи (не показано на рисунке) и тем самым обеспечивает действующую магнитодвижущую силу $\mathcal F$. На верхнем торце сердечника \textbf{(b)} зафиксирован крепёж с проводящим мостом \textbf{(g)} и двумя клеммами \textbf{(f)}. На расстоянии $x_c=\qty{3}{\mm}$ от клемм \textbf{(f)} расположены клеммы \textbf{(h)}, которые подключены к внешнему выключателю (не показано на рисунке). Если клеммы \textbf{(f)} коснутся клемм \textbf{(h)}, цепь выключателя замкнётся, приводя к отключению питания дифференциального трансформатора из пункта \taskref{10}. При $\mathcal F=0$ деформации пружин \textbf{(e)} нет; силу тяжести не учитывать; краевыми эффектами пренебречь.

\begin{task}
    \item При каком минимальном $\mathcal F=\mathcal F_0$ магнитный пускатель замкнётся и дифференциальный трансформатор отключится от сети питания? Если всего $N=1000$ витков \textbf{(d)}, чему при каком численном значении пороговой амплитудной силы тока на обмотке магнитного пускателя $I_M$ сработает предохранитель? Сравните результат с \taskref{10}. Рассматривайте квазистатическую динамику процесса.
\end{task}

Эта задача показывает, что дифференциальный трансформатор является отличным устройством защитного отключения --- небольшие изменения подаваемого в трансформатор тока ввиду внешней утечки (через человека, влагу, землю, и т.д.), которые приводят к $I_1\neq I_2$ в пункте \taskref{10}, приводит к срабатыванию защитного отключения!

\cpic{0.48}{figures/3.6.png}[Магнитный пускатель к пункту \taskref{11}]